% ============================================================
\section{Limitations and Research Gaps}
% ============================================================

An honest characterisation of the system's current boundaries is necessary for
researchers designing studies around it.

\begin{description}
	\item[Yaw drift in prolonged ambulation.]
	      The magnetometer-free design eliminates magnetic interference but accepts a slow
	      yaw drift during extended free ambulation.
	      For tasks that require absolute yaw accuracy (e.g., three-dimensional rotation
	      analysis during rotation-heavy sport), the kinematic correction algorithms may be
	      insufficient without periodic reset manoeuvres.

	\item[Ballistic movement transients.]
	      The current filtering stack is optimised for movements below approximately
	      \SI{60}{\degree\per\second}.
	      High-speed ballistic actions (golf swing, sprint, jump landing) produce \acs{IMU}
	      transients that inflate short-term position error.
	      Sport-specific validation studies are needed before deploying the system in
	      these contexts.

	\item[Upper cervical coverage.]
	      The longest sensor variant (45\,cm) covers the sacrum through the upper thoracic
	      spine.
	      Cervical kinematics require a separate cervical attachment, which is not
	      commercially available at present.

	\item[Masch et al.\ validation data.]
	      The geometric accuracy figures cited in this document are drawn from an
	      in-vitro study referenced in Walkling et al.\ \cite{walkling2025} but not yet
	      published as a standalone peer-reviewed paper.
	      Independent replication and full methodological reporting of this validation remain
	      outstanding.

	\item[Population diversity in \ac{HAR} models.]
	      The Walkling et al.\ HAR study involved 30 participants \cite{walkling2025}.
	      Generalisation to populations with pathological spinal morphology (scoliosis,
	      post-surgical fusion, severe osteoporosis) has not been assessed.

	\item[Long-term sensor drift and mechanical fatigue.]
	      Printed strain gauges may exhibit electrical resistance drift and mechanical
	      hysteresis after extended use and repeated laundering cycles.
	      Long-term test data on measurement stability across sensor lifetime have not been
	      published.
\end{description}
